\clearpage
\chapter*{DAFTAR SIMBOL}
\addcontentsline{toc}{chapter}{DAFTAR SIMBOL}

\noindent\fbox{%
\begin{minipage}{\textwidth}
\textit{\textbf{Pedoman Penulisan Daftar Simbol:} Daftar ini berisi simbol-simbol ilmiah, notasi matematika, variabel, singkatan teknis, atau parameter khusus yang digunakan dalam laporan. Simbol-simbol tersebut harus didefinisikan secara jelas beserta satuannya (jika ada) untuk memudahkan pembaca dalam memahami isi tulisan.}
\end{minipage}%
}
\vspace{0.5cm}

\begin{longtable}{ll}
\textbf{Simbol} & \textbf{Keterangan} \\
\hline
$\alpha$ & Konstanta pembelajaran (\textit{learning rate}) pada jaringan saraf tiruan \\
$\sum$ & Simbol penjumlahan matematika (\textit{summation}) \\
$\sigma$ & Standar deviasi dari sebaran data penelitian \\
$G$ & Representasi Graph (simpul dan sisi) pada pemodelan jaringan \\
$w(u,v)$ & Bobot jarak dari simpul $u$ ke simpul $v$ \\
$Q$ & Antrian (\textit{Queue}) simpul yang belum dioptimasi \\
\hline
\end{longtable}
